\begin{abstract}
Most memory systems for chat assistants organize what they store as the conversation unfolds, merging and summarizing
facts before any question is asked. Whether this organization helps is unclear, because published scores confound the memory with the reader model, the answer prompt, the context budget, and the judge. We hold
these fixed on BEAM, whose conversations reach a million tokens, and find that the answer prompt and the context
budget alone open gaps as large as those between systems. We then test the alternative: \method{} stores atomic
records once and never rewrites them, and organizes them only after the question arrives, compiling retrieved records
into notes on time order, changes, conflicts, missing coverage, and the user's standing instructions. With gpt-6-luna
as the reader, it beats Hindsight, a strong write-time system given the same extra call, by \prelim{7.3} and \prelim{4.2} points on the 100K and 1M tiers (mean of three runs) while using fewer read-time tokens in total; with three stronger readers it is ahead by 1--3
points in every case, resolved in one. The lesson is to store what was said and organize it for the question; in our experiments, the benefit shrinks as the reader gets stronger. All results are on development sets, and gpt-6-luna serves as reader, memory builder, and judge, so the largest gains may partly reflect a judge preference. \todo{held-out test results}
\end{abstract}
